\documentclass[12pt,a4paper]{ctexart}
\usepackage[utf8]{inputenc}
\usepackage{geometry}
\geometry{a4paper, margin=1in}
\usepackage{hyperref}
\usepackage{graphicx}
\usepackage{tcolorbox}
\usepackage{xcolor}
\usepackage{listings}
\usepackage{setspace}
\usepackage{fancyhdr}
\usepackage{titlesec}
\usepackage{tocloft}

\tcbuselibrary{skins,breakable}

% 段落间距与行高
\setlength{\parskip}{0.8em}
\setstretch{1.4}

% 标题颜色定制
\titleformat{\section}
  {\normalfont\Large\bfseries\color{blue!80!black}}
  {\thesection}{1em}{}
  
\titleformat{\subsection}
  {\normalfont\large\bfseries\color{teal!80!black}}
  {\thesubsection}{1em}{}

% 设置页面样式
\pagestyle{fancy}
\fancyhf{}
\fancyhead[L]{\textbf{零基础网络排错指南}}
\fancyhead[R]{\thepage}
\renewcommand{\headrulewidth}{0.4pt}
\setlength{\headheight}{15pt}

% 自定义颜色
\definecolor{mygreen}{rgb}{0,0.6,0}
\definecolor{mygray}{rgb}{0.5,0.5,0.5}
\definecolor{mymauve}{rgb}{0.58,0,0.82}
\definecolor{logbg}{rgb}{0.15,0.15,0.15}
\definecolor{logfg}{rgb}{0.9,0.9,0.9}

% 链接颜色
\hypersetup{
    colorlinks=true,
    linkcolor=blue!70!black,
    filecolor=magenta,      
    urlcolor=cyan,
}

% 配置代码块样式 (针对一般代码)
\lstset{ 
  backgroundcolor=\color{white},   
  basicstyle=\ttfamily\small,        
  breaklines=true,                 
  captionpos=b,                    
  commentstyle=\color{mygreen},    
  keywordstyle=\color{blue},       
  stringstyle=\color{mymauve},     
  frame=single,
  rulecolor=\color{black!30},
}

% 配置终端日志样式
\lstdefinestyle{terminal}{
  backgroundcolor=\color{logbg},
  basicstyle=\ttfamily\small\color{logfg},
  breaklines=true,
  frame=none,
  moredelim=[s][\color{green!80!white}]{*}{\ },
  moredelim=[s][\color{red!80!white}]{<}{\ },
  moredelim=[s][\color{yellow!80!white}]{>}{\ }
}

\begin{document}

% 自定义紧凑型封面
\begin{center}
    \vspace*{0.5cm}
    {\Huge \textbf{\color{blue!80!black}计算机网络与代理排错全书}}\\[0.8cm]
    {\LARGE \textbf{结合终端实战的零基础指南}}\\[0.5cm]
    {\large 作者：写给小白的通俗指南}\\[0.2cm]
    {\small (包含了从0到1的底层逻辑、日志分析、双网卡冲突与代码级解释)}\\[0.5cm]
    {\normalsize \today}
    \vspace*{0.8cm}
    \hrule height 1pt
    \vspace*{0.5cm}
\end{center}

% 目录紧凑排版
\renewcommand{\contentsname}{\centering \Large 目录}
\setlength{\cftbeforesecskip}{0.3em}
\tableofcontents

\vspace{1cm}
\hrule height 1pt
\vspace{1cm}

% ================= 第一章 =================
\section{引言：把网络世界想象成“寄快递”}

如果你没有计算机基础，不要害怕那些复杂的英文单词。整个计算机网络，其实就是一个\textbf{巨大的“寄快递”和“收发信件”}的系统。

\begin{tcolorbox}[colback=blue!5,colframe=blue!75!black,title=概念对照表,fonttitle=\bfseries,boxrule=0.5pt,arc=3mm]
\begin{itemize}
    \item \textbf{你想看一个网页} $\rightarrow$ 你要给别人寄一封信，要求对方把网页数据给你寄回来。
    \item \textbf{数据包 (Packet)} $\rightarrow$ 就是你寄出去的那封信。
    \item \textbf{路由器 (Router)} $\rightarrow$ 就是各个城市的顺丰/京东快递中转站。
    \item \textbf{IP 地址 (IP Address)} $\rightarrow$ 就是门牌号（你的发件地址和目标的收件地址）。
    \item \textbf{端口号 (Port)} $\rightarrow$ 某栋大楼里的具体房间号，或者是具体的收件人。
    \item \textbf{代理 (Proxy)} $\rightarrow$ 强制设立在你家门口的“快递代收发站”。
\end{itemize}
\end{tcolorbox}

本讲义将完全基于你电脑终端中发生的一起真实的“包裹丢失（网络不通）”惨案，带你彻底揭开计算机网络的底层秘密。

\newpage
% ================= 第二章 =================
\section{第一部分：地址的秘密（IP 与 CIDR）}

要送信，首先得知道门牌号。但在计算机世界里，门牌号分为很多种类型。

\subsection{1. 什么是 IP 地址？}
IP 地址是设备在网络中的身份证号。在我们这次的案例中，出现了两个极其特殊的门牌号：
\begin{itemize}
    \item \textbf{127.0.0.1（回环地址）}：这个地址永远代表\textbf{“我自己”}。就像你左手递东西给右手，根本不需要出门。电脑发往这个地址的信件，会在操作系统的内部直接“折返”，不需要真正的网卡和网线。这也是为什么即使你拔掉网线，你依然可以访问 127.0.0.1。
    \item \textbf{192.168.x.x（局域网 IP）}：这是你家\textbf{“小区内部”}的门牌号。比如你家的大模型服务地址是 \texttt{192.168.5.9}，你手机可能是 \texttt{192.168.5.10}。在小区里面（连接同一个WiFi或路由器）串门非常快。但是，小区外面的人（公网）是绝对找不到这个地址的。
\end{itemize}

\subsection{2. 什么是 CIDR（比如 /16）？}
在网络配置和日志中，你经常会看到类似 \texttt{192.168.0.0/16} 这样的写法。
这里的 \texttt{/16} 究竟是什么意思？为什么它能把系统搞崩溃？

CIDR 的全称是 \textbf{Classless Inter-Domain Routing（无类别域间路由）}。你可以把它理解为\textbf{“邮政编码”或者“小区统称”}。
它规定了：“只要门牌号前两段是 192.168 的，都是同一个小区的！”

\begin{tcolorbox}[colback=yellow!10,colframe=orange!80!black,title=硬核知识：掩码的二进制计算,fonttitle=\bfseries,boxrule=0.5pt,arc=3mm]
在计算机底层，CIDR 是如何判断 \texttt{192.168.5.9} 是不是属于 \texttt{192.168.0.0/16} 的？
\begin{enumerate}
    \item \textbf{转为二进制：} 计算机把 IP 切成 32 个 0 和 1。
    \item \textbf{/16 的含义：} 意思是“前16个数字必须严格匹配，后面的不用管”。
    \item 换算下来，\texttt{/16} 对应的是子网掩码 \texttt{255.255.0.0}。
    \item 计算机将目标IP与掩码进行“按位与（AND）”运算，发现匹配成功。它就知道：“哦，这是自己人！”
\end{enumerate}
\end{tcolorbox}

\newpage
% ================= 第三章 =================
\section{第二部分：快递代收点（Proxy 与环境变量）}

\subsection{1. 什么是 Proxy (代理)？}
Proxy 在英文里是“代理人”的意思。你可以把它当成一个\textbf{“快递代收发站”}。
当你没有代理时，你的电脑（快递员）直接出门送信。当有了代理（比如 Clash 软件），你的电脑会把信交给 Clash，Clash 帮你寄给远方，收到回信再交给你。

\subsection{2. 系统怎么知道有代收发站？（环境变量）}
操作系统为了告诉所有的软件“去哪里找这个代收发站”，设立了\textbf{环境变量}。这相当于在你家的大门上贴了两张显眼的告示：

\begin{tcolorbox}[colback=red!5,colframe=red!75!black,title=第一张告示：强制代收规则 (http\_proxy),boxrule=0.5pt,arc=3mm]
\texttt{export http\_proxy="http://127.0.0.1:7890"} \\
\textbf{大白话：}“所有寄出去的信，一律先交给 7890 号代收发站（Clash）处理！不管目的地是哪！”
\end{tcolorbox}

\begin{tcolorbox}[colback=green!5,colframe=green!75!black,title=第二张告示：豁免白名单 (no\_proxy),boxrule=0.5pt,arc=3mm]
\texttt{export no\_proxy="192.168.0.0/16"} \\
\textbf{大白话：}“注意：如果要寄给小区邻居（/16 邮编范围），就自己亲自去送，**不要**交给代收发站！”
\end{tcolorbox}

了解了门牌号和门口的这两张告示，我们马上进入案发现场。

\newpage
% ================= 第四章 =================
\section{第三部分：案发现场日志剖析（为什么信寄丢了？）}

根据你提供的真实终端日志，我们将案件分为四幕进行彻底剖析。

\subsection{第一幕：发出去的信，石沉大海}
\textbf{【原始日志分析】}
\begin{lstlisting}[style=terminal]
● Bash(curl -s http://192.168.5.9:3000/ 2>&1 | head -50)
  __  (No output)
\end{lstlisting}

\textbf{【发生了什么？】}
你雇佣了一个名叫“\textbf{curl}”的命令行快递员，对他说：“去敲局域网门牌号为 \texttt{192.168.5.9} 的 3000 号房间，把网页数据取回来。”
结果，快递员去了，空着手回来了。终端没有任何输出内容。

\subsection{第二幕：给快递员装上“行车记录仪”}
为了找出问题，你给快递员加上了 \texttt{-v}（verbose 详细模式）参数，这相当于开启了他的行车记录仪。

\textbf{【原始日志分析】}
\begin{lstlisting}[style=terminal]
● Bash(curl -v http://192.168.5.9:3000/ 2>&1 | head -40)
* Uses proxy env variable no_proxy == 'localhost,127.0.0.1,192.168.0.0/16...'
* Uses proxy env variable http_proxy == 'http://127.0.0.1:7890/'
*   Trying 127.0.0.1:7890...
* Connected to (nil) (127.0.0.1) port 7890 (#0)
> GET http://192.168.5.9:3000/ HTTP/1.1
> Host: 192.168.5.9:3000
...
< HTTP/1.1 502 Bad Gateway
\end{lstlisting}

\newpage
\textbf{【日志逐行破解与底层逻辑】}
\begin{itemize}
    \item \texttt{Uses proxy env variable...}: 快递员出门前，确实看到了你贴在门上的两张告示（读取了环境变量）。
    \item \texttt{Trying 127.0.0.1:7890...}: 令人震惊的一幕出现了！快递员居然无视了白名单，把原本应该在内网递交的信件，死死地塞给了本地代理端口（Clash代收发站）。这叫做\textbf{“应用层的自愿流量劫持”}。
    \item \texttt{< HTTP/1.1 502 Bad Gateway}: 最终收到了退件！为什么是 502？
\end{itemize}

\begin{tcolorbox}[colback=purple!5,colframe=purple!75!black,title=深度原理：502 错误究竟是谁报的？,fonttitle=\bfseries,boxrule=0.5pt,arc=3mm]
很多人以为 502 是目标网站挂了。错！
在这个场景下，\textbf{502 是 Clash 代收发站给你盖的退件戳}。
代收发站拿着这封信，跑去国外的中转站找。可是外网根本找不到你家小区的私有地址 \texttt{192.168.5.9}。于是代收发站只能把信退回来，告诉你：**502 网关错误（我没坏，但上家连不上，查无此路）**。
\end{tcolorbox}

\subsection{第三幕：破案！死板的快递员与 CIDR 盲区}
真正的悬念来了：你明明在告示里写了白名单 \texttt{192.168.0.0/16}，为什么信还是交给了代收发站？

\textbf{【原始日志分析】}
\begin{lstlisting}[style=terminal]
Bash(curl --version 2>&1 | head -1)
curl 7.81.0 (x86_64-pc-linux-gnu)...
\end{lstlisting}

原因出在快递员身上。你雇佣的老版快递员（\texttt{curl 7.81.0}）\textbf{文化程度太低}。

\begin{tcolorbox}[colback=red!10,colframe=red!80!black,title=版本代沟造成的“弱智”行为,fonttitle=\bfseries,boxrule=0.5pt,arc=3mm]
\texttt{curl 7.81.0} 的底层 C 语言代码中，并没有实现 CIDR 的二进制按位运算功能。
他把目标地址 \texttt{192.168.5.9} 和纸条上的 \texttt{192.168.0.0/16} 做了个“大家找不同”。
他一看：“哎哟，这两个词的字母长得完全不一样啊！带有斜杠的什么鬼？”
于是他死板地判定：**不在白名单内，交代收发站！**
\end{tcolorbox}

\newpage
\subsection{大结局：下达死命令，强夺控制权}
为了拿到数据，你使用了非常暴力的覆盖指令。

\textbf{【原始日志分析】}
\begin{lstlisting}[style=terminal]
● Bash(curl -s --noproxy '*' http://192.168.5.9:3000/ )
<!doctype html>...
\end{lstlisting}

\textbf{【技术解析】}
你在命令行中增加的 \texttt{--noproxy '*'} 参数，优先级是全宇宙最高的。它强制覆盖了系统门口贴的告示。
你揪着快递员的耳朵下死命令：\textbf“把之前看的告示全忘了！不管纸条上写什么，你给我亲自走到 192.168.5.9 去！”
这一次，快递员终于直接通过 ARP 协议在局域网内找到了目标，拿回了网页数据。

% ================= 第五章 =================
\section{第四部分：遇到问题怎么解决（给小白的三招）}

以后遇到网络连不通、代理报错，不要慌，按照以下三个层级去排查和解决。

\subsection{第一招：查监控（永远加 \texttt{-v}）}
这是每个程序员必备的条件反射。遇到 \texttt{curl} 报错，立刻加上 \texttt{-v} 参数。
\begin{lstlisting}[language=bash]
curl -v http://你的目标网站.com
\end{lstlisting}
如果输出里有 \texttt{* Trying 127.0.0.1:7890...}，说明流量被代理劫持了。

\subsection{第二招：防坑测试（\texttt{--noproxy '*'}）}
当你怀疑代理把网络搞乱了，想证明“如果没有代理，我是不是就能连上？”，请使用这个命令：
\begin{lstlisting}[language=bash]
curl --noproxy '*' http://你的目标网站.com
\end{lstlisting}
如果加上这句立马能连上，那么 100\% 是代理配置（比如白名单失效）的问题。

\subsection{第三招：终极解法，告别死板规则}
\begin{enumerate}
    \item \textbf{治本法（解雇笨员工）：} 把系统的 \texttt{curl} 升级到 \texttt{7.86.0} 以上版本，新版原生支持 CIDR 解析。
    \item \textbf{治标法（把话说直白）：} 把白名单里的 \texttt{/16} 去掉，直接写死目标门牌号：\texttt{export no\_proxy="192.168.5.9"}。
    \item \textbf{降维打击法（开启 TUN 模式）：} 在 Clash 中开启 TUN（虚拟网卡）模式。这就相当于修了自动传送带。开启 TUN 后，不用再贴环境变量告示，只要数据一出门就会掉进系统的虚拟地道，由 Clash 底层通过更高级的 Trie 树算法自动帮你完美分拣。
\end{enumerate}

\newpage
% ================= 第六章 =================
\section{第五部分：双路并行——插上网线开WiFi，电脑傻了怎么办？}

很多同学会遇到一个极其经典的工业现场问题：
\textbf{“我电脑开着 WiFi（为了上网），同时插了一根网线连着工业机械臂（IP 是 192.168.5.1）。结果我必须把 WiFi 关掉，才能连上机械臂。这是为什么？怎么解决？”}

要解开这个谜团，我们要引入计算机网络中极其重要的两个概念：“网卡”和“路由表”。

\subsection{1. 两扇门与门卫的烦恼（什么是网卡与路由表？）}
\textbf{网卡（Network Interface）：} 
你的电脑只要连了网，就会开一扇门。连着 WiFi，就是开了一扇“无线门”；插着网线，就是开了一扇“有线门”。
现在，你的电脑同时开了两扇门。

\textbf{路由表（Routing Table）：}
电脑内部有一个名叫“路由表”的门卫老大。每当你要寄一封信时，门卫老大就要决定：\textbf{“这封信，该从哪扇门丢出去？”}

\subsection{2. 为什么开着 WiFi 门卫就会“懵逼”？}

分析你电脑底层的真实路由表：
\begin{lstlisting}[style=terminal]
$ ip addr
2: enp3s0: <BROADCAST,MULTICAST,UP,LOWER_UP> ... state UP
    inet6 fe80::d693:90ff:fe3d:c5b8/64 scope link 
3: wlp0s20f3: <BROADCAST,MULTICAST,UP,LOWER_UP> ...
    inet 192.168.8.9/24 ...

$ ip route get 192.168.5.1
192.168.5.1 via 198.18.0.2 dev Meta table 2022 src 198.18.0.1
\end{lstlisting}

\begin{tcolorbox}[colback=red!5,colframe=red!75!black,title=案发现场：门卫与 Clash 的连环套,fonttitle=\bfseries,boxrule=0.5pt,arc=3mm]
通过你电脑的真实数据，我们发现了两个致命秘密：
\begin{enumerate}
    \item \textbf{有线网卡（enp3s0）根本没有 IPv4 地址！} 只有一串长长的 IPv6 地址。这意味着你的电脑在“有线门”上，连个合法的门牌号都没有！
    \item \textbf{流量被 Clash TUN 彻底劫持！} 当你尝试访问机械臂（192.168.5.1）时，路由表显示它通过 \texttt{dev Meta} 发送。而 \texttt{Meta} 正是 Clash 开启的虚拟网卡！
\end{enumerate}
\end{tcolorbox}

\textbf{为什么必须关闭 WiFi 才能连上？}
因为你的有线网卡没有配 IP，电脑认为自己和 \texttt{192.168.5.x} 网段“不熟”。此时只要开着 WiFi，所有的流量（包括找机械臂的信）都会无脑掉进 Clash 的虚拟地道（TUN 接口 \texttt{Meta}）里，被送往外网，机械臂当然不会有任何响应。
当你关闭 WiFi 时，Clash 传送带和 WiFi 大门全部关闭，电脑在无路可走的情况下，才可能会尝试通过有线网卡进行最原始的广播寻找。

\subsection{3. 终极解决办法（小白必看步骤）}

要解决这个问题，你根本不需要频繁开关 WiFi，只需要给你的有线网卡**配置一个静态 IP**。

打开你电脑的“网络适配器设置（选择有线网卡 enp3s0 的 IPv4 设置）”，手动填写以下参数：
\begin{itemize}
    \item \textbf{IP 地址：} 填入与机械臂同小区的任意地址，比如 \texttt{192.168.5.10}。
    \item \textbf{子网掩码：} \texttt{255.255.255.0}（代表 /24 小区）。
    \item \textbf{默认网关：} \textbf{打死也不要填！留空！} 
    \item \textbf{DNS 服务器：} \textbf{留空！}
\end{itemize}

\begin{tcolorbox}[colback=green!5,colframe=green!75!black,title=为什么这样配就能解决？,fonttitle=\bfseries,boxrule=0.5pt,arc=3mm]
一旦你给有线网卡配置了 \texttt{192.168.5.10}，电脑的路由表里就会瞬间生成一条高优先级的局部路由规则：
\textbf{“凡是去往 192.168.5.x 小区的信，一律直接从有线网卡（enp3s0）扔出去！”} \\
这条局部规则的优先级（掩码长度 /24）远远高于 Clash 传送带的默认全局规则（/0）。
从此，即便你开着 WiFi、挂着 Clash，只要你访问 \texttt{192.168.5.1}，门卫老大也会雷厉风行地一脚把信踢进有线网线里，直达机械臂！而你的上网流量依然会快乐地走 WiFi 门。
\end{tcolorbox}

\newpage
% ================= 第七章 =================
\section{第六部分：番外篇——为什么 TUN 模式能“降维打击”？}

你刚才亲自测试了开启 Clash 的 TUN 模式，发现原来那个死活连不上的问题瞬间消失了。这究竟是施了什么魔法？什么是“Trie 树算法”？

\subsection{1. 为什么 TUN 能解决旧版软件的 Bug？}

我们要回到之前的比喻：
\textbf{不开 TUN 时（环境变量代理）}：是快递员（curl）出门前自己看门上的告示（\texttt{no\_proxy}）。因为快递员太笨看不懂 \texttt{/16} 的小区编码，所以他送错了。这就叫“应用层代理”——极其依赖软件自身的智商。

\textbf{开启 TUN 时（虚拟网卡代理）}：
你把大门上的告示全撕了（清空环境变量）。快递员出门时啥也不用想，拿着信直接往前走。
但是！你在大门口挖了一个坑，装了一部\textbf{隐形的传送带}（TUN 虚拟网卡）。
不管快递员多笨，他一迈出门槛，信件就会自动掉进这个传送带里，被直接送到 Clash 的“地下秘密分拣中心”。

到了分拣中心，由极其聪明的 Clash 站长亲自来看这封信。Clash 站长是懂 \texttt{192.168.0.0/16} 这种 CIDR 编码的。他一看信封上写着 \texttt{192.168.5.9}，立刻大手一挥：“这是同一个小区的信，走直连通道（Direct），不用发给国外的代收发站！”
于是，信件顺利送达。

\begin{tcolorbox}[colback=purple!5,colframe=purple!75!black,title=硬核知识：OSI 模型的降维打击,fonttitle=\bfseries,boxrule=0.5pt,arc=3mm]
环境变量代理工作在 \textbf{OSI 第七层（应用层）}，每个应用（curl, python, 浏览器）都有自己的实现逻辑，很容易出 Bug。
TUN 模式工作在 \textbf{OSI 第三层（网络层）}。它通过修改操作系统的底层路由表，强制接管了所有通过网卡发出的 IP 数据包。应用层根本不知道自己被劫持了。这就是所谓的“降维打击”。
\end{tcolorbox}

\subsection{2. 什么是 Trie 树（前缀树）算法？}

现在，全世界有成千上万个小区编码（CIDR 规则），Clash 站长每天要处理几万封信，他为什么能在一秒钟内判断出这封信该走国内还是国外？这就靠 \textbf{Trie 树（前缀树）}。

\textbf{如果不用 Trie 树（笨办法）：}
站长拿着手里的小本本，从第一页开始一行一行对：
“192.168.5.9 是 10.0.0.0 吗？不是。”
“是 172.16.0.0 吗？不是。”
这样找下去，电脑早就卡死了。

\textbf{如果用了 Trie 树（字典树算法）：}
Trie 树就像是\textbf{查字典}。
站长看到门牌号的第一位是 \texttt{192}，他直接翻到字典的 \texttt{192} 这一章；
下一位是 \texttt{168}，他顺着目录直接翻到 \texttt{168} 这一页；
然后是 \texttt{5}，他一秒钟就定位到了 \texttt{192.168.5.x} 这个抽屉，发现上面贴着标签：\textbf{直连（Direct）}！

\begin{tcolorbox}[colback=blue!5,colframe=blue!75!black,title=Trie 树的威力,fonttitle=\bfseries,boxrule=0.5pt,arc=3mm]
通过 Trie 树（前缀匹配），无论 Clash 里装了 10 条规则还是 10 万条规则，查找的速度几乎是一模一样的！它只需要顺着 IP 地址的每一位往下找（对于 IPv4 最多找 32 次），这在计算机里只要零点零几毫秒。这就是 Clash 能做到“完美且无感分拣”的底层性能密码。
\end{tcolorbox}

\newpage
% ================= 第七章 =================
\section{第六部分：DNS——为什么必须要有个“接线员”？}

在刚才修改路由器的配置中，有一个名为 \textbf{DNS 服务器} 的选项。很多初学者总是把它和网关搞混。
什么是 DNS (Domain Name System)？在“寄快递”的比喻里，它扮演了什么角色？

\subsection{1. 互联网的超级“114查号台”}

我们在上网时，输入的是 \texttt{www.baidu.com} 或者 \texttt{github.com}。
但是！快递员（电脑）只认得\textbf{纯数字的门牌号}（比如 \texttt{39.156.66.10}）。你如果直接给快递员一个名字（比如“马化腾的家”），快递员是没法送信的，因为地图上根本没有画“名字”，只有“经纬度坐标”。

这时候，就需要 \textbf{DNS（超级查号台 / 电话簿）}了。

\begin{tcolorbox}[colback=purple!5,colframe=purple!75!black,title=DNS 的工作流程,fonttitle=\bfseries,boxrule=0.5pt,arc=3mm]
\textbf{你：} “快递员，帮我去 \texttt{www.baidu.com} 拿个东西。” \\
\textbf{快递员（电脑）：} “我不认识这几个英文字母啊！等一下，我先打个电话给 114 查号台（也就是配置里的 DNS 服务器）。” \\
\textbf{快递员：} “喂？114吗？请问 \texttt{www.baidu.com} 的真实物理门牌号是多少？” \\
\textbf{DNS服务器（查号台）：} “稍等，我查一下小本本... 查到了，它的门牌号是 \texttt{39.156.66.10}。” \\
\textbf{快递员：} “收到！”（然后快递员才真正骑上车，朝着 \texttt{39.156.66.10} 绝尘而去）。
\end{tcolorbox}

\subsection{2. 为什么路由器里要把 DNS 设为 192.168.10.1？}

在局域网（你家）里，你的 WiFi 路由器（\texttt{192.168.10.1}）不仅是掌控大门的“门卫老大（网关）”，它还**兼职**做“片区 114 查号接线员”。
当你的电脑想要查询 \texttt{www.baidu.com} 时，会直接问路由器。路由器内部如果不知道，路由器会再代替你去问“国家级的总查号台（比如电信、联通的 DNS，或者 114.114.114.114）”。

\textbf{如果不填 DNS 会怎样？}
如果你的电脑没有配置 DNS，你打开浏览器输入 \texttt{www.baidu.com} 会瞬间报错：\textbf{无法解析服务器的 DNS 地址}。
因为快递员手里拿着收件人的名字，却死活找不到能查门牌号的电话簿，最后只能原地罢工。
有趣的是，如果你当时直接输入的是数字 IP（比如 \texttt{192.168.5.9} 或者 \texttt{39.156.66.10}），那么不配 DNS 也一样能连上，因为快递员根本不需要去查电话簿了！

\subsection{3. 代理（Clash）环境下的 DNS 终极魔法：防污染}

很多时候，即使你设了正规的查号台，你依然上不去某些国外的网站（比如 Google）。这不是因为 Google 坏了，而是因为\textbf{“查号台（DNS）被人动了手脚”}。
国内的查号台在遇到你查询 Google 时，会故意给你报一个**假的门牌号**（比如告诉你去太平洋中心），这叫做 \textbf{DNS 污染}。

\begin{tcolorbox}[colback=red!5,colframe=red!75!black,title=Clash 的 Fake-IP（假查号）黑科技,fonttitle=\bfseries,boxrule=0.5pt,arc=3mm]
为了解决查号台骗人的问题，Clash 使出了一招极其不可思议的绝技：\textbf{发假驾照}。
当你的电脑向查号台查询 \texttt{google.com} 时，Clash 瞬间把这通电话掐断！然后 Clash 自己伪装成查号台，随便给你发一个假的门牌号（比如 \texttt{198.18.0.1}）。\\
快递员拿着这个假门牌号刚出门，立刻就掉进了我们前面说的 \textbf{TUN 虚拟传送带} 里！ \\
Clash 在地下室截获了这个假门牌号，笑着说：“哈哈，这假门牌号是我发的，我知道你其实想去 Google。来，这封信我亲自通过加密通道发到国外的真代理服务器，让国外的代理服务器帮你查真正的门牌号！” \\
这就完美避开了国内查号台的欺骗。
\end{tcolorbox}

\newpage
% ================= 第八章 =================
\section{附录一：进阶极客实验 —— 手写 Python 模拟 CIDR 计算}

如果你觉得光听“按位与”运算不过瘾，我们可以用 10 行 Python 代码，亲自模拟一下电脑底层是如何判断门牌号是否在同一个小区的。

\begin{lstlisting}[language=Python, title=cidr\_check.py]
import ipaddress

# 1. 定义你的小区白名单 (CIDR)
my_neighborhood = ipaddress.ip_network("192.168.0.0/16")

# 2. 拿到这封信的目标门牌号
target_ip = ipaddress.ip_address("192.168.5.9")

# 3. 让计算机做底层比对
if target_ip in my_neighborhood:
    print(f"[{target_ip}] 是自己人！不要交给代收发站，亲自去送！")
else:
    print(f"[{target_ip}] 是外人，把它交给代收发站！")
\end{lstlisting}

当你运行这段代码时，输出结果一定是“自己人”。
而在老版本的 curl 源代码（C语言）中，它缺乏上面这种聪明的库，只能做：
\begin{lstlisting}[language=C, title=笨笨的C语言逻辑]
if (strcmp("192.168.5.9", "192.168.0.0/16") == 0) {
    // 只有长得完全一样才放行
    deliver_directly();
} else {
    send_to_proxy();
}
\end{lstlisting}

这就是我们常说的：**“技术架构上的降维打击”**。

\newpage
% ================= 第九章 =================
\section{附录二：TCP 三次握手的生活化比喻}

在快递员出发之前，还需要经过一个“确认接头”的环节，也就是传说中的 TCP 三次握手。小白听到这个词往往觉得高深莫测，其实它就像是在打电话。

\subsection{第一步：SYN（你好，听得到吗？）}
快递员（curl）到达目标地址 \texttt{192.168.5.9} 门外，他并没有直接把包裹扔进去。
他先在门口按响了门铃，并大喊一声：“喂！3000 号房间有人吗？我有一个快递要交接，你能听到我说话吗？”
这一声呼喊，在网络协议里叫做 \textbf{SYN 包}。

\subsection{第二步：SYN-ACK（听得到！我也准备好了）}
3000 号房间的主人（服务器程序）听到门铃后，走到门边回答：
“听得到！我这边已经准备好接收包裹了。那你能听到我说话吗？”
这句回应，在网络协议里叫做 \textbf{SYN-ACK 包}。

\subsection{第三步：ACK（好的，那我们开始吧！）}
快递员听到主人的回应，确认了对方是个活人并且能正常交流，于是最后回答一句：
“好的，我也听到你了，我们正式开始交接包裹吧！”
这最后一句确认，在网络协议里叫做 \textbf{ACK 包}。

\begin{tcolorbox}[colback=green!5,colframe=green!75!black,title=为什么非要三次？,fonttitle=\bfseries,boxrule=0.5pt,arc=3mm]
很多人问，为什么不能两次？如果只有两次：
快递员喊：“有人吗？”
主人答：“有！”
然后主人就开始傻等。万一主人回答“有”的这句话被风吹走了（网络丢包），快递员根本没听见，快递员就走了。留下主人一个人在房间里一直等，浪费了感情和时间。
只有经过**三次**确认，双方才能 100\% 确定对方既能听见自己，也能向自己说话。这是一条极为严谨的通信逻辑。
\end{tcolorbox}

\vspace{2cm}
\begin{center}
\textbf{\huge 全书完}
\end{center}

\end{document}